% --- CORRECCION: Forzar interpretacion raw de la entrada ---
\UseRawInputEncoding
\documentclass[aspectratio=169, 10pt, xcolor=table]{beamer}

% --- Configuracion del Tema ---
\usetheme{Madrid}
\usecolortheme{dolphin}

% --- Paquetes necesarios ---
\usepackage[utf8]{inputenc}
\usepackage[T1]{fontenc}
\usepackage[spanish,es-noquoting]{babel}
\usepackage{amsmath, amssymb}
\usepackage{mathtools}
\usepackage{booktabs}
\usepackage{graphicx}
\usepackage{listings}
\usepackage{tikz}
\usetikzlibrary{arrows.meta, positioning, shapes.geometric, babel}
\usepackage{etoolbox}
\usepackage{upquote}

% --- Estilo para codigo SQL ---
\definecolor{codegreen}{rgb}{0,0.6,0}
\definecolor{codegray}{rgb}{0.5,0.5,0.5}
\definecolor{codepurple}{rgb}{0.58,0,0.82}
\definecolor{backcolour}{rgb}{0.95,0.95,0.92}
\lstset{
    language=SQL,
    backgroundcolor=\color{backcolour},
    commentstyle=\color{codegreen},
    keywordstyle=\color{blue},
    numberstyle=\tiny\color{codegray},
    stringstyle=\color{codepurple},
    basicstyle=\ttfamily\scriptsize,
    breaklines=true,
    frame=single,
    showstringspaces=false,
    literate={á}{{\'a}}1 {é}{{\'e}}1 {í}{{\'i}}1 {ó}{{\'o}}1 {ú}{{\'u}}1 {ñ}{{\~n}}1 {ü}{{\"u}}1
}

% --- Pie de pagina personalizado ---
\makeatletter
\setbeamertemplate{footline}{
  \leavevmode%
  \hbox{%
  \begin{beamercolorbox}[wd=.333333\paperwidth,ht=2.25ex,dp=1ex,center]{author in head/foot}%
    \usebeamerfont{author in head/foot}\insertshortauthor
  \end{beamercolorbox}%
  \begin{beamercolorbox}[wd=.333333\paperwidth,ht=2.25ex,dp=1ex,center]{title in head/foot}%
    \usebeamerfont{title in head/foot}Sesión 7
  \end{beamercolorbox}%
  \begin{beamercolorbox}[wd=.333333\paperwidth,ht=2.25ex,dp=1ex,right]{date in head/foot}%
    \usebeamerfont{date in head/foot}BBDD \hfill \insertframenumber/\inserttotalframenumber \hspace*{0.3cm}
  \end{beamercolorbox}}%
  \vskip0pt%
}
\makeatother

% --- Datos de la portada ---
\title[SESIÓN 7 BBDD]{\textbf{Sesión 7}}
\subtitle{Programación en el Servidor y Objetos de Base de Datos}
\author[Prof. C. Sánchez]{Carlos César Sánchez Coronel}
\institute{\textbf{BASES DE DATOS}}
\date{}

\setbeamertemplate{navigation symbols}{}

\AtBeginSection[]{
  \begin{frame}
    \frametitle{Contenido de la sección}
    \tableofcontents[currentsection]
  \end{frame}
}

\begin{document}

\begin{frame}
    \titlepage
\end{frame}

\begin{frame}{Contenido general}
    \tableofcontents
\end{frame}

% ============================================================================
% SECCIÓN 1: INTRODUCCIÓN
% ============================================================================
\section{Introducción}

\begin{frame}{Objetivos de la sesión}
    \begin{itemize}
        \item Comprender el papel de los objetos de base de datos: vistas, funciones, procedimientos, triggers.
        \item Aprender a crear y utilizar estos objetos para encapsular lógica de negocio.
        \item Conocer técnicas de hash y cifrado para datos sensibles.
        \item Integrar la base de datos con aplicaciones externas, ETL y modelos de IA.
    \end{itemize}
\end{frame}

% ============================================================================
% SECCIÓN 2: VISTAS
% ============================================================================
\section{Vistas}

\begin{frame}[fragile]{Vistas Simples}
    \begin{block}{Definición}
        Tabla virtual basada en una consulta. No almacena datos físicamente.
    \end{block}
    \begin{exampleblock}{Ejemplo}
        \begin{lstlisting}
CREATE VIEW clientes_activos AS
SELECT id, nombre, email
FROM clientes
WHERE activo = true;

SELECT * FROM clientes_activos WHERE ciudad = 'Madrid';
        \end{lstlisting}
    \end{exampleblock}
    \begin{itemize}
        \item \textbf{Seguridad}: ocultar columnas sensibles.
        \item \textbf{Simplicidad}: encapsular joins complejos.
    \end{itemize}
\end{frame}

\begin{frame}[fragile]{Vistas Materializadas}
    \begin{block}{Características}
        Almacenan físicamente el resultado de la consulta. Ideales para reportes.
    \end{block}
    \begin{exampleblock}{PostgreSQL}
        \begin{lstlisting}
CREATE MATERIALIZED VIEW resumen_ventas AS
SELECT producto_id, SUM(cantidad) AS total_vendido
FROM ventas GROUP BY producto_id;

REFRESH MATERIALIZED VIEW resumen_ventas;
REFRESH MATERIALIZED VIEW CONCURRENTLY resumen_ventas;
        \end{lstlisting}
    \end{exampleblock}
    \begin{itemize}
        \item \textbf{Ventaja}: consultas rápidas.
        \item \textbf{Desventaja}: datos no están en tiempo real.
    \end{itemize}
\end{frame}

\begin{frame}[fragile]{Ejercicio Resuelto: Vista para Informes}
    \textbf{Enunciado:} Crear una vista que muestre para cada cliente su nombre, total gastado y número de pedidos (solo clientes con al menos un pedido).
    \begin{block}{Solución}
        \begin{lstlisting}
CREATE VIEW resumen_clientes AS
SELECT c.id, c.nombre,
       COUNT(p.id) AS num_pedidos,
       COALESCE(SUM(p.total), 0) AS total_gastado
FROM clientes c
LEFT JOIN pedidos p ON c.id = p.cliente_id
GROUP BY c.id, c.nombre
HAVING COUNT(p.id) > 0;
        \end{lstlisting}
    \end{block}
\end{frame}

% ============================================================================
% SECCIÓN 3: FUNCIONES Y PROCEDIMIENTOS
% ============================================================================
\section{Funciones y Procedimientos Almacenados}

\begin{frame}{Diferencias Clave}
    \begin{block}{Función}
        \begin{itemize}
            \item Retorna un valor (escalar o tabla).
            \item Se usa dentro de expresiones SQL.
            \item No puede realizar transacciones (commit/rollback) en la mayoría de motores.
        \end{itemize}
    \end{block}
    \begin{block}{Procedimiento}
        \begin{itemize}
            \item Puede no retornar valor o devolver múltiples valores.
            \item Puede incluir transacciones.
            \item Se invoca con \texttt{CALL}.
        \end{itemize}
    \end{block}
\end{frame}

\begin{frame}[fragile]{Funciones Escalares}
    \begin{exampleblock}{Ejemplo: Calcular IVA (PostgreSQL)}
        \begin{lstlisting}
CREATE OR REPLACE FUNCTION calcular_iva(monto NUMERIC)
RETURNS NUMERIC AS $$
BEGIN
    RETURN COALESCE(monto, 0) * 1.21;
END;
$$ LANGUAGE plpgsql IMMUTABLE;

SELECT calcular_iva(100); -- 121
        \end{lstlisting}
    \end{exampleblock}
    \begin{exampleblock}{Ejemplo: MySQL}
        \begin{lstlisting}
CREATE FUNCTION calcular_iva(monto DECIMAL(10,2))
RETURNS DECIMAL(10,2) DETERMINISTIC
BEGIN
    RETURN monto * 1.21;
END;
        \end{lstlisting}
    \end{exampleblock}
\end{frame}

\begin{frame}[fragile]{Funciones Tabla}
    \begin{block}{Retornan un conjunto de filas}
        \begin{lstlisting}
CREATE OR REPLACE FUNCTION empleados_por_depto(depto_id INT)
RETURNS TABLE(id INT, nombre TEXT, salario NUMERIC) AS $$
BEGIN
    RETURN QUERY
    SELECT id, nombre, salario
    FROM empleados
    WHERE departamento_id = depto_id;
END;
$$ LANGUAGE plpgsql STABLE;

SELECT * FROM empleados_por_depto(2);
        \end{lstlisting}
    \end{block}
\end{frame}

\begin{frame}[fragile]{Procedimientos Almacenados}
    \begin{block}{Ejemplo: Transferencia bancaria (PostgreSQL)}
        \begin{lstlisting}
CREATE OR REPLACE PROCEDURE transferir(
    origen INT, destino INT, monto NUMERIC
)
LANGUAGE plpgsql AS $$
BEGIN
    UPDATE cuentas SET saldo = saldo - monto WHERE id = origen;
    UPDATE cuentas SET saldo = saldo + monto WHERE id = destino;
    COMMIT;
END;
$$;

CALL transferir(1, 2, 100);
        \end{lstlisting}
    \end{block}
\end{frame}

\begin{frame}[fragile]{Caso de Uso Real: Generación de Tickets en Gasolinera}
    \textbf{Problema:} Generar ticket correlativo por estación y fecha, con concurrencia.
    \begin{block}{Solución con función}
        \begin{lstlisting}
CREATE TABLE ventas_gasolina (
    id SERIAL PRIMARY KEY,
    estacion_id INT NOT NULL,
    fecha DATE NOT NULL DEFAULT CURRENT_DATE,
    ticket_num INT NOT NULL,
    importe NUMERIC NOT NULL,
    UNIQUE (estacion_id, fecha, ticket_num)
);

CREATE OR REPLACE FUNCTION generar_ticket(p_estacion INT, p_importe NUMERIC)
RETURNS INT AS $$
DECLARE
    v_ticket INT;
BEGIN
    SELECT COALESCE(MAX(ticket_num), 0) + 1 INTO v_ticket
    FROM ventas_gasolina
    WHERE estacion_id = p_estacion AND fecha = CURRENT_DATE
    FOR UPDATE;

    INSERT INTO ventas_gasolina (estacion_id, fecha, ticket_num, importe)
    VALUES (p_estacion, CURRENT_DATE, v_ticket, p_importe);

    RETURN v_ticket;
END;
$$ LANGUAGE plpgsql;
        \end{lstlisting}
    \end{block}
\end{frame}

% ============================================================================
% SECCIÓN 4: TRIGGERS
% ============================================================================
\section{Triggers}

\begin{frame}{Tipos de Triggers}
    \begin{description}
        \item[BEFORE] Se ejecuta antes de la operación. Útil para validaciones o modificar valores.
        \item[AFTER] Después de la operación. Ideal para auditoría, sincronización.
        \item[INSTEAD OF] Reemplaza la operación (en vistas).
    \end{description}
    \vspace{0.3cm}
    \begin{description}
        \item[FOR EACH ROW] Se ejecuta por cada fila afectada.
        \item[FOR EACH STATEMENT] Se ejecuta una vez por sentencia.
    \end{description}
\end{frame}

\begin{frame}[fragile]{Trigger de Auditoría}
    \begin{block}{Registrar cambios de salario}
        \begin{lstlisting}
CREATE TABLE audit_salarios (
    id_empleado INT,
    salario_anterior NUMERIC,
    salario_nuevo NUMERIC,
    fecha TIMESTAMP DEFAULT CURRENT_TIMESTAMP,
    usuario TEXT DEFAULT current_user
);

CREATE OR REPLACE FUNCTION audit_salario_func()
RETURNS TRIGGER AS $$
BEGIN
    IF OLD.salario IS DISTINCT FROM NEW.salario THEN
        INSERT INTO audit_salarios (id_empleado, salario_anterior, salario_nuevo)
        VALUES (OLD.id, OLD.salario, NEW.salario);
    END IF;
    RETURN NEW;
END;
$$ LANGUAGE plpgsql;

CREATE TRIGGER trig_audit_salario
AFTER UPDATE OF salario ON empleados
FOR EACH ROW EXECUTE FUNCTION audit_salario_func();
        \end{lstlisting}
    \end{block}
\end{frame}

\begin{frame}[fragile]{Trigger de Validación}
    \begin{block}{Evitar salarios negativos}
        \begin{lstlisting}
CREATE OR REPLACE FUNCTION validar_salario()
RETURNS TRIGGER AS $$
BEGIN
    IF NEW.salario < 0 THEN
        RAISE EXCEPTION 'El salario no puede ser negativo';
    END IF;
    RETURN NEW;
END;
$$ LANGUAGE plpgsql;

CREATE TRIGGER trig_validar_salario
BEFORE INSERT OR UPDATE ON empleados
FOR EACH ROW EXECUTE FUNCTION validar_salario();
        \end{lstlisting}
    \end{block}
\end{frame}

\begin{frame}{Precauciones con Triggers}
    \begin{itemize}
        \item \textbf{Rendimiento}: overhead en cada operación DML, especialmente por fila.
        \item \textbf{Recursividad}: pueden dispararse a sí mismos; controlar con condiciones.
        \item \textbf{Depuración}: lógica oculta, difícil de rastrear.
        \item \textbf{Mantenibilidad}: efectos laterales no evidentes.
    \end{itemize}
\end{frame}

% ============================================================================
% SECCIÓN 5: PROGRAMACIÓN AVANZADA
% ============================================================================
\section{Programación Avanzada}

\begin{frame}[fragile]{Cursores}
    \begin{block}{Recorrer filas una por una (ineficiente, pero a veces necesario)}
        \begin{lstlisting}
CREATE OR REPLACE FUNCTION procesar_empleados()
RETURNS VOID AS $$
DECLARE
    cur CURSOR FOR SELECT id, salario FROM empleados WHERE activo = true;
    v_id INT;
    v_salario NUMERIC;
BEGIN
    OPEN cur;
    LOOP
        FETCH cur INTO v_id, v_salario;
        EXIT WHEN NOT FOUND;
        UPDATE empleados SET salario = v_salario * 1.05 WHERE id = v_id;
    END LOOP;
    CLOSE cur;
END;
$$ LANGUAGE plpgsql;
        \end{lstlisting}
    \end{block}
\end{frame}

\begin{frame}[fragile]{Manejo de Excepciones}
    \begin{block}{Capturar errores en procedimientos}
        \begin{lstlisting}
CREATE OR REPLACE PROCEDURE transferir_seguro(
    origen INT, destino INT, monto NUMERIC
)
LANGUAGE plpgsql AS $$
BEGIN
    UPDATE cuentas SET saldo = saldo - monto WHERE id = origen;
    UPDATE cuentas SET saldo = saldo + monto WHERE id = destino;
    COMMIT;
EXCEPTION
    WHEN OTHERS THEN
        ROLLBACK;
        RAISE NOTICE 'Error en transferencia: %', SQLERRM;
END;
$$;
        \end{lstlisting}
    \end{block}
\end{frame}

% ============================================================================
% SECCIÓN 6: HASH Y CRIPTOGRAFÍA
% ============================================================================
\section{Funciones de Hash y Criptografía}

\begin{frame}{Conceptos Fundamentales}
    \begin{description}
        \item[Hash] Función unidireccional (no reversible). Ej: MD5, SHA-256.
        \item[Cifrado] Reversible mediante clave (simétrico/asimétrico).
        \item[Salt] Valor aleatorio añadido al hash para evitar ataques de diccionario.
    \end{description}
\end{frame}

\begin{frame}[fragile]{Funciones Hash en SQL}
    \begin{block}{PostgreSQL (requiere pgcrypto)}
        \begin{lstlisting}
CREATE EXTENSION IF NOT EXISTS pgcrypto;
SELECT encode(digest('mi_secreto', 'sha256'), 'hex');
        \end{lstlisting}
    \end{block}
    \begin{block}{MySQL}
        \begin{lstlisting}
SELECT MD5('mi_secreto');
SELECT SHA2('mi_secreto', 256);
        \end{lstlisting}
    \end{block}
    \begin{block}{SQL Server}
        \begin{lstlisting}
SELECT HASHBYTES('SHA2_256', 'mi_secreto');
        \end{lstlisting}
    \end{block}
\end{frame}

\begin{frame}[fragile]{Almacenamiento Seguro de Contraseñas}
    \begin{block}{Insertar usuario con hash + salt (PostgreSQL)}
        \begin{lstlisting}
INSERT INTO usuarios (email, password_hash, salt)
VALUES (
    'usuario@mail.com',
    encode(digest('password123' || gen_salt('md5'), 'sha256'), 'hex'),
    gen_salt('md5')
);
        \end{lstlisting}
    \end{block}
    \begin{block}{Verificación}
        \begin{lstlisting}
SELECT * FROM usuarios
WHERE email = 'usuario@mail.com'
AND password_hash = encode(digest('password_ingresada' || salt, 'sha256'), 'hex');
        \end{lstlisting}
    \end{block}
\end{frame}

\begin{frame}[fragile]{Anonimización de Datos}
    \begin{block}{Sustituir emails por hash (entornos de test)}
        \begin{lstlisting}
UPDATE clientes
SET email = encode(digest(email || 's3cr3t', 'sha256'), 'hex');
        \end{lstlisting}
    \end{block}
    \begin{alertblock}{Precaución}
        Salt fijo no es seguro si el conjunto de emails es pequeño. Usar salt por registro.
    \end{alertblock}
\end{frame}

% ============================================================================
% SECCIÓN 7: INTEGRACIÓN CON SISTEMAS EXTERNOS
% ============================================================================
\section{Integración con Sistemas Externos}

\begin{frame}{¿Por qué integrar lógica en la base de datos?}
    \begin{itemize}
        \item \textbf{Latencia}: código cerca de los datos.
        \item \textbf{Consistencia}: reglas centralizadas.
        \item \textbf{Seguridad}: control de acceso a nivel de fila/columna.
        \item \textbf{Rendimiento}: operaciones complejas optimizadas por el motor.
    \end{itemize}
\end{frame}

\begin{frame}[fragile]{Llamada desde Python (psycopg2)}
    \begin{block}{Ejemplo con función y procedimiento}
        \begin{lstlisting}[language=Python]
import psycopg2

conn = psycopg2.connect("dbname=test user=postgres")
cur = conn.cursor()

# Llamar a función que retorna valor
cur.callproc('generar_ticket', [5, 45.50])
ticket = cur.fetchone()[0]
print(f"Ticket: {ticket}")

# Llamar a procedimiento
cur.callproc('transferir', [1, 2, 100])
conn.commit()
cur.close()
conn.close()
        \end{lstlisting}
    \end{block}
\end{frame}

\begin{frame}[fragile]{Integración con Machine Learning}
    \begin{block}{Función que expone datos para entrenamiento}
        \begin{lstlisting}
CREATE OR REPLACE FUNCTION datos_entrenamiento(fecha_inicio DATE, fecha_fin DATE)
RETURNS TABLE(fecha DATE, producto_id INT, total_vendido NUMERIC) AS $$
BEGIN
    RETURN QUERY
    SELECT fecha, producto_id, total_vendido
    FROM resumen_diario
    WHERE fecha BETWEEN fecha_inicio AND fecha_fin;
END;
$$ LANGUAGE plpgsql STABLE;
        \end{lstlisting}
    \end{block}
    \begin{block}{Consumo en Python}
        \begin{lstlisting}[language=Python]
import pandas as pd
df = pd.read_sql("SELECT * FROM datos_entrenamiento('2025-01-01', '2025-03-01')", conn)
        \end{lstlisting}
    \end{block}
\end{frame}

\begin{frame}[fragile]{Orquestación con Apache Airflow}
    \begin{block}{DAG que ejecuta un procedimiento diario}
        \begin{lstlisting}[language=Python]
from airflow import DAG
from airflow.providers.postgres.operators.postgres import PostgresOperator
from datetime import datetime

with DAG('etl_diario', start_date=datetime(2025,1,1), schedule='@daily') as dag:
    tarea = PostgresOperator(
        task_id='actualizar_resumen',
        postgres_conn_id='postgres_default',
        sql="CALL actualizar_resumen_diario();"
    )
        \end{lstlisting}
    \end{block}
\end{frame}

\begin{frame}{Seguridad: SECURITY DEFINER vs INVOKER}
    \begin{description}
        \item[SECURITY INVOKER] (por defecto) Se ejecuta con privilegios del que llama.
        \item[SECURITY DEFINER] Se ejecuta con privilegios del creador.
    \end{description}
    \begin{exampleblock}{Ejemplo}
        \begin{lstlisting}
CREATE OR REPLACE FUNCTION eliminar_usuario(uid INT)
RETURNS VOID SECURITY DEFINER AS $$
BEGIN
    DELETE FROM usuarios WHERE id = uid;
END;
$$ LANGUAGE plpgsql;

REVOKE ALL ON usuarios FROM app_user;
GRANT EXECUTE ON FUNCTION eliminar_usuario(INT) TO app_user;
        \end{lstlisting}
    \end{exampleblock}
\end{frame}

% ============================================================================
% SECCIÓN 8: COMPARATIVA ENTRE MOTORES
% ============================================================================
\section{Comparativa entre Motores}

\begin{frame}{Características de Programación}
    \begin{table}
        \centering
        \begin{tabular}{l|c|c|c}
            \toprule
            Característica & PostgreSQL & MySQL & SQL Server \\
            \midrule
            Lenguaje procedural & PL/pgSQL, Python, etc. & SQL/PSM & T-SQL \\
            Funciones escalares & Sí & Sí & Sí \\
            Funciones tabla & Sí (RETURNS TABLE) & Sí (RETURNS) & Sí (TVF) \\
            Procedimientos & Sí (CREATE PROCEDURE) & Sí & Sí \\
            Triggers & BEFORE/AFTER & BEFORE/AFTER & INSTEAD OF, AFTER \\
            Cursores & Sí & Sí & Sí \\
            Vistas materializadas & Sí (REFRESH) & No nativas & Indexed views \\
            SECURITY DEFINER & Sí & SQL SECURITY & EXECUTE AS \\
            \bottomrule
        \end{tabular}
    \end{table}
\end{frame}

% ============================================================================
% SECCIÓN 9: OPTIMIZACIÓN DE CÓDIGO
% ============================================================================
\section{Optimización de Código}

\begin{frame}{Buenas Prácticas}
    \begin{itemize}
        \item Usar operaciones basadas en conjuntos en lugar de cursores.
        \item Evitar triggers innecesarios.
        \item Indexar columnas usadas en condiciones dentro de funciones.
        \item Marcar funciones como \texttt{IMMUTABLE}, \texttt{STABLE} o \texttt{VOLATILE} según corresponda.
        \item Analizar consultas con \texttt{EXPLAIN}.
    \end{itemize}
\end{frame}

\begin{frame}[fragile]{Ejemplo de Optimización}
    \begin{block}{Ineficiente (cursor)}
        \begin{lstlisting}
FOR r IN SELECT salario FROM empleados WHERE departamento_id = p_depto LOOP
    total := total + r.salario;
END LOOP;
        \end{lstlisting}
    \end{block}
    \begin{block}{Eficiente (SQL directo)}
        \begin{lstlisting}
SELECT SUM(salario) INTO total FROM empleados WHERE departamento_id = p_depto;
        \end{lstlisting}
    \end{block}
\end{frame}

% ============================================================================
% SECCIÓN 10: EJERCICIOS RESUELTOS INTEGRADORES
% ============================================================================
\section{Ejercicios Resueltos}

\begin{frame}[fragile]{Ejercicio 1: Trigger de Stock}
    \textbf{Enunciado:} Al insertar en \texttt{detalle\_pedido}, actualizar stock y abortar si insuficiente.
    \begin{block}{Solución}
        \begin{lstlisting}
CREATE OR REPLACE FUNCTION actualizar_stock()
RETURNS TRIGGER AS $$
DECLARE
    v_stock INT;
BEGIN
    SELECT stock INTO v_stock FROM productos WHERE id = NEW.producto_id FOR UPDATE;
    IF v_stock < NEW.cantidad THEN
        RAISE EXCEPTION 'Stock insuficiente para producto %', NEW.producto_id;
    END IF;
    UPDATE productos SET stock = stock - NEW.cantidad WHERE id = NEW.producto_id;
    RETURN NEW;
END;
$$ LANGUAGE plpgsql;

CREATE TRIGGER trig_actualizar_stock
BEFORE INSERT ON detalle_pedido
FOR EACH ROW EXECUTE FUNCTION actualizar_stock();
        \end{lstlisting}
    \end{block}
\end{frame}

\begin{frame}[fragile]{Ejercicio 2: Función para Calcular Total de Pedido}
    \textbf{Enunciado:} Dado un \texttt{pedido\_id}, calcular la suma de (cantidad * precio) y actualizar la columna \texttt{total} en \texttt{pedidos}.
    \begin{block}{Solución}
        \begin{lstlisting}
CREATE OR REPLACE FUNCTION calcular_total_pedido(p_pedido_id INT)
RETURNS NUMERIC AS $$
DECLARE
    v_total NUMERIC;
BEGIN
    SELECT SUM(dp.cantidad * p.precio) INTO v_total
    FROM detalle_pedido dp
    JOIN productos p ON dp.producto_id = p.id
    WHERE dp.pedido_id = p_pedido_id;

    UPDATE pedidos SET total = v_total WHERE id = p_pedido_id;
    RETURN v_total;
END;
$$ LANGUAGE plpgsql;
        \end{lstlisting}
    \end{block}
\end{frame}

\begin{frame}[fragile]{Ejercicio 3: Vista con Datos Sensibles Ocultos}
    \textbf{Enunciado:} Crear vista que muestre nombre y salario, pero oculte salario a usuarios no autorizados.
    \begin{block}{Solución}
        \begin{lstlisting}
CREATE VIEW empleados_rh AS
SELECT id, nombre,
       CASE WHEN CURRENT_USER = 'rh_user' THEN salario ELSE NULL END AS salario,
       departamento_id
FROM empleados;
        \end{lstlisting}
    \end{block}
    Luego otorgar permisos select a los usuarios.
\end{frame}

\begin{frame}[fragile]{Ejercicio 4: Transferencia con Auditoría}
    \textbf{Enunciado:} Procedimiento que transfiera dinero entre cuentas y registre en auditoría.
    \begin{block}{Solución}
        \begin{lstlisting}
CREATE TABLE auditoria_transferencias (
    id SERIAL PRIMARY KEY,
    origen INT, destino INT, monto NUMERIC,
    fecha TIMESTAMP DEFAULT CURRENT_TIMESTAMP,
    usuario TEXT DEFAULT current_user
);

CREATE OR REPLACE PROCEDURE transferir_con_auditoria(
    p_origen INT, p_destino INT, p_monto NUMERIC
)
LANGUAGE plpgsql AS $$
DECLARE
    v_saldo NUMERIC;
BEGIN
    SELECT saldo INTO v_saldo FROM cuentas WHERE id = p_origen FOR UPDATE;
    IF v_saldo < p_monto THEN
        RAISE EXCEPTION 'Saldo insuficiente';
    END IF;

    UPDATE cuentas SET saldo = saldo - p_monto WHERE id = p_origen;
    UPDATE cuentas SET saldo = saldo + p_monto WHERE id = p_destino;
    INSERT INTO auditoria_transferencias (origen, destino, monto)
    VALUES (p_origen, p_destino, p_monto);
    COMMIT;
END;
$$;
        \end{lstlisting}
    \end{block}
\end{frame}

% ============================================================================
% SECCIÓN 11: GLOSARIO
% ============================================================================
\section{Glosario}

\begin{frame}{Términos Clave}
    \begin{description}
        \item[Vista] Tabla virtual basada en una consulta.
        \item[Vista materializada] Vista que almacena físicamente los datos.
        \item[Función escalar] Retorna un valor simple.
        \item[Función tabla] Retorna un conjunto de filas.
        \item[Procedimiento almacenado] Bloque de código que realiza acciones.
        \item[Trigger] Código que se ejecuta automáticamente ante un evento.
        \item[BEFORE / AFTER / INSTEAD OF] Momento de ejecución del trigger.
        \item[FOR EACH ROW / STATEMENT] Nivel del trigger.
        \item[Hash] Función unidireccional (resumen).
        \item[Salt] Valor aleatorio añadido al hash.
        \item[SECURITY DEFINER / INVOKER] Modo de ejecución con permisos.
    \end{description}
\end{frame}

% --- Diapositiva final ---
\begin{frame}
    \centering

    
    \vspace{1cm}
    \large ¡Gracias!
\end{frame}

\end{document}